% Part: first-order-logic
% Chapter: natural-deduction
% Section: soundness

\documentclass[../../../include/open-logic-section]{subfiles}

\begin{document}

\iftag{FOL}
      {\olfileid{fol}{ntd}{sou}}
      {\olfileid{pl}{ntd}{sou}}
      
\olsection{Correctheid}

\begin{explain}
!!^a{derivation}ssysteem, zoals natuurlijke deductie, is
\emph{correct} als het alleen dingen kan !!{derive} die werkelijk
volgen. Correctheid geeft dus een veiligheidsgarantie voor
!!{derivation}ssystemen. Welke precieze eigenschap we nodig hebben,
hangt van de situatie af. We willen bijvoorbeeld weten dat
\begin{enumerate}
\item elke !!{derivable} !!{sentence}
  \iftag{FOL}{geldig}{onder elke valuatie waar} is;
\item als !!a{sentence} !!{derivable} is uit andere zinnen, zij daar
  ook logisch uit volgt;
\item als een verzameling !!{sentence}s inconsistent is, zij niet waar
  te maken is.
\end{enumerate}
Dit zijn belangrijke eigenschappen van !!a{derivation}ssysteem. Als
een ervan niet geldt, kan het systeem te veel !!{derive} en deugt het
niet. Daarom is het heel belangrijk om de correctheid van
!!a{derivation}ssysteem te bewijzen.
\end{explain}

\begin{thm}[Correctheid]
\ollabel{thm:soundness}
Als $!A$ !!{derivable} is uit de !!{undischarged} aannamen
$\Gamma$, dan $\Gamma \Entails !A$.
\end{thm}

\begin{proof}
Laat $\delta$ !!a{derivation} van $!A$ zijn. We gebruiken inductie
over het aantal inferenties in~$\delta$.

In de inductiebasis is het aantal inferenties~$0$. Dan bestaat
$\delta$ alleen uit één !!{sentence}~$!A$: een aanname. Die aanname is
!!{undischarged}. Alleen een inferentie kan een aanname
!!{discharged} maken, en hier zijn geen inferenties. Elke
\iftag{FOL}{!!{structure}~$\Struct{M}$}{!!{valuation}~$\pAssign{v}$}
die alle !!{undischarged} aannamen van het bewijs waar maakt, maakt
daarom ook~$!A$ waar.

Nu komt de inductiestap. Neem aan dat $\delta$~$n$ inferenties bevat.
De uitgangsformule of uitgangsformules van de onderste inferentie
worden !!{derive}d met deel-!!{derivation}s. Elk daarvan heeft minder
dan~$n$ inferenties. De inductiehypothese zegt dat de uitgangsformules
van de onderste inferentie volgen uit de !!{undischarged} aannamen van
de deel-!!{derivation}s die daar eindigen. We moeten laten zien dat de
conclusie~$!A$ volgt uit de !!{undischarged} aannamen van het hele
bewijs.

We behandelen elk mogelijk soort onderste inferentie apart. Eerst
komen de inferenties met maar één uitgangsformule.

\begin{enumerate}
\item Stel dat de laatste inferentie \Intro{\lnot} is. De
  !!{derivation} heeft dan deze vorm:
  \begin{prooftree}
    \AxiomC{$\Gamma, \Discharge{!A}{n}$}
    \RightLabel{$\delta_1$}
    \DeduceC{$\lfalse$}
    \DischargeRule{\Intro{\lnot}}{n}
    \UnaryInfC{$\lnot !A$}
  \end{prooftree}
  Volgens de inductiehypothese volgt $\lfalse$ uit de
  !!{undischarged} aannamen $\Gamma \cup \{!A\}$ van~$\delta_1$.
  Beschouw
\iftag{FOL}{!!a{structure}~$\Struct{M}$}{!!a{valuation}~$\pAssign{v}$}.
  We moeten laten zien: als
  $\iftag{FOL}{\Sat{M}}{\pSat{v}}{\Gamma}$, dan
  $\iftag{FOL}{\Sat{M}}{\pSat{v}}{\lnot !A}$. Neem voor een bewijs via
  tegenspraak aan dat $\iftag{FOL}{\Sat{M}}{\pSat{v}}{\Gamma}$, maar
  $\iftag{FOL}{\Sat/{M}}{\pSat/{v}}{\lnot !A}$. Dit laatste betekent
  $\iftag{FOL}{\Sat{M}}{\pSat{v}}{!A}$. Dan zou
  $\iftag{FOL}{\Sat{M}}{\pSat{v}}{\Gamma \cup \{!A\}}$ gelden. Dat is
  in strijd met de inductiehypothese. Dus
  $\iftag{FOL}{\Sat{M}}{\pSat{v}}{\lnot !A}$.

\item De laatste inferentie is \Elim{\land}. Er zijn twee
  mogelijkheden: we mogen $!A$ of $!B$ uit de uitgangsformule
  $!A \land !B$ afleiden. Bekijk de eerste mogelijkheid. De
  !!{derivation}~$\delta$ ziet er zo uit:
  \begin{prooftree}
    \AxiomC{$\Gamma$}
    \RightLabel{$\delta_1$}
    \DeduceC{$!A \land !B$}
    \RightLabel{\Elim{\land}}
    \UnaryInfC{$!A$}
  \end{prooftree}
  Volgens de inductiehypothese volgt $!A \land !B$ uit de
  !!{undischarged} aannamen~$\Gamma$ van~$\delta_1$. Beschouw
  !!a{structure}~\iftag{FOL}{$\Struct{M}$}{$\pAssign{v}$}. We moeten
  laten zien: als $\iftag{FOL}{\Sat{M}}{\pSat{v}}{\Gamma}$, dan
  $\iftag{FOL}{\Sat{M}}{\pSat{v}}{!A}$. Neem aan dat
  $\iftag{FOL}{\Sat{M}}{\pSat{v}}{\Gamma}$. De inductiehypothese
  ($\Gamma \Entails !A \land !B$) geeft
  $\iftag{FOL}{\Sat{M}}{\pSat{v}}{!A \land !B}$. Per definitie geldt
  $\iftag{FOL}{\Sat{M}}{\pSat{v}}{!A \land !B}$ precies als zowel
  $\iftag{FOL}{\Sat{M}}{\pSat{v}}{!A}$ als
  $\iftag{FOL}{\Sat{M}}{\pSat{v}}{!B}$. (De mogelijkheid waarin $!B$
  uit $!A \land !B$ wordt afgeleid, werkt op dezelfde manier.)
  
\item De laatste inferentie is \Intro{\lor}. Er zijn twee
  mogelijkheden: we mogen $!A \lor !B$ afleiden uit de
  uitgangsformule~$!A$ of uit~$!B$. Bekijk de eerste mogelijkheid. De
  !!{derivation} heeft deze vorm:
  \begin{prooftree}
    \AxiomC{$\Gamma$}
    \RightLabel{$\delta_1$}
    \DeduceC{$!A$}
    \RightLabel{\Intro{\lor}}
    \UnaryInfC{$!A \lor !B$}
  \end{prooftree}
  Volgens de inductiehypothese volgt $!A$ uit de !!{undischarged}
  aannamen~$\Gamma$ van~$\delta_1$. Beschouw
\iftag{FOL}{!!a{structure}~$\Struct{M}$}{!!a{valuation}~$\pAssign{v}$}.
  We moeten laten zien: als
  $\iftag{FOL}{\Sat{M}}{\pSat{v}}{\Gamma}$, dan
  $\iftag{FOL}{\Sat{M}}{\pSat{v}}{!A \lor !B}$. Neem aan dat
  $\iftag{FOL}{\Sat{M}}{\pSat{v}}{\Gamma}$. Dan geldt
  $\iftag{FOL}{\Sat{M}}{\pSat{v}}{!A}$ wegens
  $\Gamma \Entails !A$ (de inductiehypothese). Daarom geldt ook
  $\iftag{FOL}{\Sat{M}}{\pSat{v}}{!A \lor !B}$. (De mogelijkheid
  waarin $!A \lor !B$ uit~$!B$ wordt afgeleid, werkt op dezelfde
  manier.)
  
\item De laatste inferentie is \Intro{\lif}. We leiden $!A \lif !B$
  af uit een deelbewijs met aanname~$!A$ en conclusie~$!B$. Dus:
  \begin{prooftree}
    \AxiomC{$\Gamma, \Discharge{!A}{n}$}
    \RightLabel{$\delta_1$}
    \DeduceC{$!B$}
    \DischargeRule{\Intro{\lif}}{n}
    \UnaryInfC{$!A \lif !B$}
  \end{prooftree}
  Volgens de inductiehypothese volgt $!B$ uit de !!{undischarged}
  aannamen van~$\delta_1$. Dus $\Gamma \cup \{!A\} \Entails
  !B$. Beschouw
\iftag{FOL}{!!a{structure}~$\Struct{M}$}{!!a{valuation}~$\pAssign{v}$}.
  De !!{undischarged} aannamen van~$\delta$ staan alleen in $\Gamma$,
  want de laatste inferentie heft $!A$ op. We moeten dus laten zien
  dat $\Gamma \Entails !A \lif !B$. Neem voor een bewijs via
  tegenspraak aan dat er een
  \iftag{FOL}{!!{structure}~$\Struct{M}$}{!!{valuation}~$\pAssign{v}$},
  is waarvoor $\iftag{FOL}{\Sat{M}}{\pSat{v}}{\Gamma}$ maar
  $\iftag{FOL}{\Sat/{M}}{\pSat/{v}}{!A \lif !B}$. Dan
  $\iftag{FOL}{\Sat{M}}{\pSat{v}}{!A}$ en
  $\iftag{FOL}{\Sat/{M}}{\pSat/{v}}{!B}$. Volgens de hypothese volgt
  $!B$ echter uit $\Gamma \cup \{!A\}$, dus
  $\iftag{FOL}{\Sat{M}}{\pSat{v}}{!B}$. Dat is een tegenspraak. Dus
  $\Gamma \Entails !A \lif !B$.

\item De laatste inferentie is \FalseInt. Hier eindigt $\delta$ zo:
  \begin{prooftree}
    \AxiomC{$\Gamma$}
    \RightLabel{$\delta_1$}
    \DeduceC{$\lfalse$}
    \RightLabel{\FalseInt}
    \UnaryInfC{$!A$}
  \end{prooftree}
  Volgens de inductiehypothese is $\Gamma \Entails \lfalse$. We
  moeten laten zien dat $\Gamma \Entails !A$. Stel dat dit niet zo is.
  Dan is er een~$\iftag{FOL}{\Struct{M}}{\pAssign{v}}$ waarvoor
    $\iftag{FOL}{\Sat{M}}{\pSat{v}}{\Gamma}$ en
    $\iftag{FOL}{\Sat/{M}}{\pSat/{v}}{!A}$. Maar
    $\iftag{FOL}{\Sat/{M}}{\pSat/{v}}{\lfalse}$ geldt altijd. Dit zou
    dus betekenen dat $\Gamma \Entails/ \lfalse$, in strijd met de
    inductiehypothese.

\item De laatste inferentie is \FalseCl. Oefening.

\iftag{FOL}{
\item De laatste inferentie is \Intro{\lforall}. Dan heeft $\delta$
  deze vorm:
  \begin{prooftree}
    \AxiomC{$\Gamma$}
    \RightLabel{$\delta_1$}
    \DeduceC{$!A(a)$}
    \RightLabel{\Intro{\lforall}}
    \UnaryInfC{$\lforall[x][!A(x)]$}
  \end{prooftree}
  Volgens de inductiehypothese volgt de uitgangsformule $!A(a)$ uit de
  !!{undischarged} aannamen $\Gamma$. Neem een structuur
  $\Struct{M}$ die aan $\Sat{M}{\Gamma}$ voldoet. We moeten laten zien
  dat $\Sat{M}{\lforall[x][!A(x)]}$. Omdat
  $\lforall[x][!A(x)]$ !!a{sentence} is, moeten we voor elke
  variabelentoekenning~$s$ laten zien dat $\Sat{M}{!A(x)}[s]$
  (\olref[syn][ass]{prop:sat-quant}). $\Gamma$ bestaat alleen uit
  zinnen. Daarom geldt $\Sat{M}{!B}[s]$ voor elke
  $!B \in \Gamma$, volgens
  \olref[syn][sat]{defn:satisfaction}. Laat $\Struct{M'}$ gelijk zijn
  aan $\Struct{M}$, behalve dat $\Assign{a}{M'} = s(x)$. Omdat $a$
  niet in~$\Gamma$ voorkomt, geldt $\Sat{M'}{\Gamma}$ volgens
  \olref[syn][ext]{cor:extensionality-sent}. Uit
  $\Gamma \Entails !A(a)$ volgt $\Sat{M'}{!A(a)}$. Omdat $!A(a)$
  !!a{sentence} is, geldt $\Sat{M'}{!A(a)}[s]$ volgens
  \olref[syn][ass]{prop:sentence-sat-true}.
  $\Sat{M'}{!A(x)}[s]$ geldt precies als $\Sat{M'}{!A(a)}$, volgens
  \olref[syn][ext]{prop:ext-formulas} (bedenk dat $!A(a)$ gewoon
  $\Subst{!A(x)}{a}{x}$ is). Dus $\Sat{M'}{!A(x)}[s]$. Omdat $a$ niet
  voorkomt in~$!A(x)$, volgt uit
  \olref[syn][ext]{prop:extensionality} dat
  $\Sat{M}{!A(x)}[s]$. De variabelentoekenning $s$ was willekeurig.
  Dus $\Sat{M}{\lforall[x][!A(x)]}$.
  
\item De laatste inferentie is \Intro{\lexists}. Oefening.

\item De laatste inferentie is \Elim{\forall}. Oefening.
}{}
\end{enumerate}

Nu bekijken we de mogelijke inferenties met meer dan één
uitgangsformule: \Elim{\lor}, \Intro{\land},
\iftag{FOL}{\Elim{\lif} en \Elim{\lexists}}{en \Elim{\lif}}.
\begin{enumerate}

\item De laatste inferentie is \Intro{\land}. We leiden
  $!A \land !B$ af uit de uitgangsformules $!A$ en $!B$.
  $\delta$ heeft deze vorm:
  \begin{prooftree}
    \AxiomC{$\Gamma_1$}
    \RightLabel{$\delta_1$}
    \DeduceC{$!A$}
    \AxiomC{$\Gamma_2$}
    \RightLabel{$\delta_2$}
    \DeduceC{$!B$}
    \RightLabel{\Intro{\land}}
    \BinaryInfC{$!A \land !B$}
  \end{prooftree}
  Volgens de inductiehypothese volgt $!A$ uit de !!{undischarged}
  aannamen~$\Gamma_1$ van~$\delta_1$. Ook volgt $!B$ uit de
  !!{undischarged} aannamen~$\Gamma_2$ van~$\delta_2$. De
  !!{undischarged} aannamen van~$\delta$ zijn $\Gamma_1 \cup
  \Gamma_2$. Daarom moeten we laten zien dat
  $\Gamma_1 \cup \Gamma_2 \Entails !A \land !B$. Beschouw
  \iftag{FOL}{!!a{structure}~$\Struct{M}$}{!!a{valuation}~$\pAssign{v}$}
  waarvoor
  $\iftag{FOL}{\Sat{M}}{\pSat{v}}{\Gamma_1 \cup \Gamma_2}$. Uit
  $\iftag{FOL}{\Sat{M}}{\pSat{v}}{\Gamma_1}$ volgt
  $\iftag{FOL}{\Sat{M}}{\pSat{v}}{!A}$, want
  $\Gamma_1 \Entails !A$. En uit
  $\iftag{FOL}{\Sat{M}}{\pSat{v}}{\Gamma_2}$ volgt
  $\iftag{FOL}{\Sat{M}}{\pSat{v}}{!B}$, want
  $\Gamma_2 \Entails !B$. Samen geeft dit
  $\iftag{FOL}{\Sat{M}}{\pSat{v}}{!A \land !B}$.
  
\item De laatste inferentie is \Elim{\lor}. Oefening.

\item De laatste inferentie is \Elim{\lif}. We leiden $!B$ af uit de
  uitgangsformules $!A \lif !B$ en~$!A$. De !!{derivation}~$\delta$
  ziet er zo uit:
  \begin{prooftree}
    \AxiomC{$\Gamma_1$}
    \RightLabel{$\delta_1$}
    \DeduceC{$!A \lif !B$}
    \AxiomC{$\Gamma_2$}
    \RightLabel{$\delta_2$}
    \DeduceC{$!A$}
    \RightLabel{\Elim{\lif}}
    \BinaryInfC{$!B$}
  \end{prooftree}
  Volgens de inductiehypothese volgt $!A \lif !B$ uit de
  !!{undischarged} aannamen~$\Gamma_1$ van~$\delta_1$. Ook volgt $!A$
  uit de !!{undischarged} aannamen~$\Gamma_2$ van~$\delta_2$.
  Beschouw
\iftag{FOL}{!!a{structure}~$\Struct{M}$}{!!a{valuation}~$\pAssign{v}$}.
  We moeten laten zien: als
  $\iftag{FOL}{\Sat{M}}{\pSat{v}}{\Gamma_1 \cup \Gamma_2}$, dan
  $\iftag{FOL}{\Sat{M}}{\pSat{v}}{!B}$. Neem aan dat
  $\iftag{FOL}{\Sat{M}}{\pSat{v}}{\Gamma_1 \cup \Gamma_2}$. Uit
  $\Gamma_1 \Entails !A \lif !B$ volgt
  $\iftag{FOL}{\Sat{M}}{\pSat{v}}{!A \lif !B}$. Uit
  $\Gamma_2 \Entails !A$ volgt
  $\iftag{FOL}{\Sat{M}}{\pSat{v}}{!A}$. Dit betekent
  $\iftag{FOL}{\Sat{M}}{\pSat{v}}{!B}$. (Als namelijk
  $\iftag{FOL}{\Sat/{M}}{\pSat/{v}}{!B}$, dan zou uit
  $\iftag{FOL}{\Sat{M}}{\pSat{v}}{!A}$ volgen dat
  $\iftag{FOL}{\Sat/{M}}{\pSat/{v}}{!A \lif !B}$. Dat is in
  tegenspraak met
  ~$\iftag{FOL}{\Sat{M}}{\pSat{v}}{!A \lif !B}$.)

\item De laatste inferentie is \Elim{\lnot}. Oefening.
  
\tagitem{FOL}{De laatste inferentie is \Elim{\lexists}. Oefening.}{}
\end{enumerate}
\end{proof}

\tagprob{FOL}
\begin{prob}
Maak het bewijs van \olref[fol][ntd][sou]{thm:soundness} af.
\end{prob}
\tagendprob

\tagprob{notFOL}
\begin{prob}
Maak het bewijs van \olref[pl][ntd][sou]{thm:soundness} af.
\end{prob}
\tagendprob

\begin{cor}
\ollabel{cor:weak-soundness}
Als $\Proves !A$, dan is $!A$
\iftag{FOL}{geldig}{onder elke valuatie waar}.
\end{cor}

\begin{cor}
\ollabel{cor:consistency-soundness}
Als $\Gamma$ waar te maken is, dan is zij consistent.
\end{cor}

\begin{proof}
We bewijzen de tegengestelde uitspraak met omgekeerde richting, de
contrapositie. Neem aan dat $\Gamma$ niet consistent is. Dan
$\Gamma \Proves \lfalse$. Er is dus !!a{derivation} van $\lfalse$ uit
!!{undischarged} aannamen in~$\Gamma$. Volgens
\olref{thm:soundness} moet elke
\iftag{FOL}{!!{structure}~$\Struct{M}$}{!!{valuation}~$\pAssign{v}$}
die $\Gamma$ waar maakt, ook $\lfalse$ waar maken. Maar
$\iftag{FOL}{\Sat/{M}}{\pSat/{v}}{\lfalse}$ voor elke
\iftag{FOL}{!!{structure}~$\Struct{M}$}{!!{valuation}~$\pAssign{v}$},
dus geen enkele \iftag{FOL}{$\Struct{M}$}{$\pAssign{v}$} kan
$\Gamma$ waar maken. Met andere woorden: $\Gamma$ is niet waar te
maken.
\end{proof}
\end{document}
